\begin{figure}[ht]
    \begin{minipage}[t]{0.45\textwidth}
        \centering
        \[
            \begin{pmatrix}
                2 & 0  & 0  & 0  & 0 & 0 & 0 & 0 & 0 \\
                0 & -1 & 0  & 0  & 0 & 0 & 0 & 0 & 0 \\
                0 & 0  & -1 & 0  & 0 & 0 & 0 & 0 & 0 \\
                0 & 0  & 0  & -1 & 0 & 0 & 0 & 0 & 0 \\
                0 & 0  & 0  & 0  & 2 & 0 & 0 & 0 & 0 \\
                0 & 0  & 0  & 0  & 0 & -1 & 0 & 0 & 0 \\
                0 & 0  & 0  & 0  & 0 & 0 & -1 & 0 & 0 \\
                0 & 0  & 0  & 0  & 0 & 0 & 0 & -1 & 0 \\
                0 & 0  & 0  & 0  & 0 & 0 & 0 & 0 & 2 \\
            \end{pmatrix}
        \]
        \captionof{}{Matrix of the graph \( \Gamma_2 \) obtained from
            \( V_1 (m = 2)\)
        }
    \end{minipage}\hfill
    \begin{minipage}[t]{0.45\textwidth}
        \centering
        \[
            \begin{pmatrix}
                2 & 0  & 0  & 0  & 1 & 0 & 0 & 0 & 1 \\
                0 & -1 & 0  & 0  & 0 & \omega^2 & \omega & 0 & 0 \\
                0 & 0  & -1 & \omega  & 0 & 0 & 0 & \omega^2 & 0 \\
                0 & 0  & \omega^2  & -1 & 0 & 0 & 0 & \omega & 0 \\
                1 & 0  & 0  & 0  & 2 & 0 & 0 & 0 & 1 \\
                0 & \omega  & 0  & 0  & 0 & -1 & \omega^2 & 0 & 0 \\
                0 & \omega^2  & 0  & 0  & 0 & \omega & -1 & 0 & 0 \\
                0 & 0  & \omega  & \omega^2  & 0 & 0 & 0 & -1 & 0 \\
                1 & 0  & 0  & 0  & 1 & 0 & 0 & 0 & 2 \\
            \end{pmatrix}
        \]
        \captionof{}{Matrix of the graph \( \Gamma_4 \) obtained from
            \( W_{14} (m = 4)\)
        }
    \end{minipage}

    \vspace{0.5cm}

    \begin{minipage}[t]{0.45\textwidth}
        \centering
        \[
            \begin{pmatrix}
                2 & 0  & 0  & 0  & 1 & 0 & 0 & 0 & 1 \\
                0 & -1 & 0  & 0  & 0 & -1 & -1 & 0 & 0 \\
                0 & 0  & -1 & -1  & 0 & 0 & 0 & -1 & 0 \\
                0 & 0  & -1  & -1 & 0 & 0 & 0 & -1 & 0 \\
                1 & 0  & 0  & 0  & 2 & 0 & 0 & 0 & 1 \\
                0 & -1  & 0  & 0  & 0 & -1 & -1 & 0 & 0 \\
                0 & -1  & 0  & 0  & 0 & -1 & -1 & 0 & 0 \\
                0 & 0  & -1  & -1  & 0 & 0 & 0 & -1 & 0 \\
                1 & 0  & 0  & 0  & 1 & 0 & 0 & 0 & 2 \\
            \end{pmatrix}
        \]
        \captionof{}{Matrix of the graph \( \Gamma_6 \) obtained from
            \( U_2 (m = 6)\)
        }
    \end{minipage}\hfill
    \begin{minipage}[t]{0.45\textwidth}
        \centering
        \[
            \begin{pmatrix}
                2 & 0  & 0  & 0  & 3 & 0 & 0 & 0 & 3 \\
                0 & -1 & 0  & 0  & 0 & 0 & 0 & 0 & 0 \\
                0 & 0  & -1 & 0  & 0 & 0 & 0 & 0 & 0 \\
                0 & 0  & 0  & -1 & 0 & 0 & 0 & 0 & 0 \\
                3 & 0  & 0  & 0  & 2 & 0 & 0 & 0 & 3 \\
                0 & 0  & 0  & 0  & 0 & -1 & 0 & 0 & 0 \\
                0 & 0  & 0  & 0  & 0 & 0 & -1 & 0 & 0 \\
                0 & 0  & 0  & 0  & 0 & 0 & 0 & -1 & 0 \\
                3 & 0  & 0  & 0  & 3 & 0 & 0 & 0 & 2 \\
            \end{pmatrix}
        \]
        \captionof{}{Matrix of the graph \( \Gamma_8 \) obtained from \(
            \mathfrak{sl}_{3} (m = 8)\)
        }
    \end{minipage}
    \caption{Adjacency matrices, expressed in the generalised Pauli basis
        via Theorem~\ref{thm:gromada}, of the four nontrivial loop-free and
        undirected quantum graphs on \( (M_3(\mathbb{C}),
        3\operatorname{Tr}) \) up to
        isomorphism, one for each nontrivial dimension
    \( m \in \{2,4,6,8\} \).}
    \label{fig:four-matrices-m3}
\end{figure}
